\section{Discussion}\label{sec:discussion}

\subsection{Where the gain comes from}

Three pieces of evidence locate the gain. First, the control of the ablation: the same solver and budget on the nominal cost does not lower the spread (a ratio of \ablNominalMed), so the expected cost is the cause and the extra iterations are not. Second, the theory check: the first-order intensity variance induced by the aberration prior falls by a factor of \thPRatio{} from phase-fixed WGS to the robust hologram, the first-order model follows the draws of the baseline with a correlation of \thPearsonNom, and the robust hologram has \thSigmaZeroBfour{} percent nominal spread against \thSigmaZeroNom{} for the baseline, the price of the flatter response.
Third, the draws: the robust hologram has the lower spread on every draw of seed 1, which is what a lower response would produce and what a lucky tail would not.

The first-order model also says what the method does not do. It predicts a ratio of the 95th percentiles of about \thQLinRatio{} and the measurement gives \ratioMed. The difference is the second-order part of the response, which the design reaches only through the second derivatives of the cost that the sigma points match, and which grows with the aberration: at a Strehl ratio of about \strehlExp{} the rms phase offset is \phaseRms{} rad, and the first-order model of the robust hologram under-predicts its own mean squared spread by about a third
(\thPredSqBfour{} against \thMeasSqBfour). The gain is real and smaller than the first-order theory suggests. A design that reduced the second-order response directly, for example by including the quadratic terms of the intensity in the cost, is a candidate extension that this paper does not test.

\subsection{The level and the Strehl ratio}

The decomposition of the cost shows a second mechanism that Proposition 2 alone does not predict. At a target level of $0.90$ of the start efficiency, the expected cost of the baseline is lower than its nominal cost, because the aberrations lower the mean intensity of the traps (to \effDropOne{} of the start efficiency in the ensemble) toward the target. The level in effect encodes the mean intensity loss that the prior implies, and the ablation of the level has its maximum at the level
that matches the Strehl ratio of the ensemble. This is consistent with the sensitivity of the method to the pair of the design amplitude and the level: a prior of 0.03 waves per mode has a Strehl ratio near \strehlThree{} and wants a lower level, and the ablation of the amplitude at the fixed level of $0.90$ is therefore a conservative test of the amplitude and an unfair test of the level. A study that maps the pair is open. For a user, the consequence is that the level should be set from the measured Strehl ratio of
the residual, and not left at its default.

\subsection{Failure cases and negative results}

Four results are negative or partial. (i) The gain does not transfer to ensembles outside the design: it is \misDMed{} at 0.03 and \misEMed{} at 0.04 waves per mode, and \misGMed{} when modes up to Noll 21 are present. The robust hologram trades breadth for depth over the modes and amplitude it was given. (ii) The nominal spread rises: on the cases without aberration the robust hologram has 0.4 to 1.0 percent spread where the baseline is below 0.01 percent in most seeds, and its mean efficiency is \effRatioMed{} of the baseline's.
An array that needs the lowest nominal spread, and has no residual aberration, should not use it. (iii) For random displacements of the traps, the configuration of the acceptance test gives no gain, and the spread under displacement has a first-order slope term that the curvature of the spot does not explain. The explanation that the spot width alone bounds pointing robustness is incomplete in this model: with a displacement prior alone and enough steps the slope falls fivefold and the ratio reaches \poBLong{} at 0.10 bins at \poBLongEff{} of the baseline's efficiency; what remains is close to the curvature floor.
(iv) The first-order theory is met by the baseline and not by the robust hologram, as above.

\subsection{Relation to the practice of correction and calibration}

The method does not replace the Zernike correction of the large arrays~\cite{nogrette2014,ebadi2021,manetsch2025}; it assumes it. The correction removes the measured wavefront and the method designs for the residual that the correction leaves, with a prior that must come from the instrument. The residual of a real instrument is spatially varying~\cite{machu2025,christen2025}; the isotropic Gaussian prior over eight global modes used here is the simplest description, and a per-site or per-patch residual would give a prior per trap, which the expected-cost formulation accommodates (each sigma point would shift the phase of the traps of a patch) and this paper does not test.
It also does not replace camera feedback~\cite{nogrette2014,kim2019}: feedback corrects the intensities of one realization after the fact, while the expected-cost design lowers the spread of the family of realizations the experiment may present, between recalibrations. The two address different time scales and could be combined, which is a question for an experiment on hardware.

\subsection{What the result supports}

The result supports a limited statement: in a scalar Fourier model with an ideal device, for an array of 100 equal traps, and an aberration residual that is independent, Gaussian and of the amplitude that the design assumes, an expected-cost hologram started from phase-fixed WGS has a 95th-percentile spread about \ratioMed{} times lower at \effRatioMed{} of the efficiency, and a deterministic sigma-point set reaches it with half the evaluations of the 32 random draws that come close. It does not support a claim about hardware, about arrays of thousands of traps, about measured residuals, or about robustness to displacements of the traps in the acceptance configuration. The measured scope is also narrow in the choices that the ablation exposes: the design amplitude, the target level and the designed modes were set to the values of the ensemble, and the table of the ablation shows how quickly the ratio falls when any of them is set to another value. A user who adopts the method owes the same measurement for the instrument at hand: the residual of the correction, its modes and its amplitude, and a scan of the level around the Strehl ratio that the residual implies.
